\chapter{最后那段有限计算，到底在调整什么}\label{sec:controllers}
\section{只知道比例正确，还差一步}
模五的 $(1,1)$ 与 $(2,2)$ 比例相同，但列并不相同。上一讲的精确输送要求\emph{整列}同余。因此，先把比例调对以后，还得说明允许怎样改变共同倍数。

这不是要求实际数值变大。例如某段平均可以输出一个很小的共同倍数；把这个倍数从分析代表中取出后，本原列却可能变成另一个有限尺度。我们始终跟踪实际完整配置和分析代表之间的对应。

\section{一张整数矩阵，可以产生几种本原输出}
设输入列 $(a,b)$ 的最大公因数为一。施加整数矩阵后，两坐标可能有公因子。把这个最大公因数叫作\textbf{内容}，记为 $\cont$。

最简单的例子是 $A=\diag(1,6)$。输入 $(s,t)$ 本原，输出 $(s,6t)$ 的公因子是 $\gcd(s,6)$，所以只可能为 $1,2,3,6$。

甚至可以固定输入模五都等于 $(1,1)$，而让四种内容分别出现：
\[
\begin{array}{c|c|c}
\text{输入}&A\text{ 的输出}&\text{本原输出}\\ \hline
(1,1)&(1,6)&(1,6)\\
(16,1)&(16,6)&(8,3)\\
(21,1)&(21,6)&(7,2)\\
(6,1)&(6,6)&(1,1)
\end{array}
\]
它们分别除去了 $1,2,3,6$。所以同一个模五输入类，在施加同一个 $A$ 并取本原代表以后，会得到几种不同列。这就是“内容分支”的具体意思。
\InsertProof{09_terminal-unit2}
\InsertProof{lem:content}
该引理为什么要先有 $\Gamma(q,\Z)$？因为我们需要先把当前实际列送到具有指定内容的那个提升，再执行 $A$。没有这个先行准备，就不能把“某个提升存在”当成“当前输入能走过去”。

\section{有限集合上的正向操作为什么可以组成群}
假设一个置换在六步后回到原样，那么它的逆就是再做五步。这是有限群中最朴素的事实。在本题中，每条安全内容分支是有限标记集上的置换，所以无需额外要求它的有理矩阵具有统一逆词，仍能在这个有限集合里反向移动。

更妙的是，先走一次内容为 $g$ 的分支，再走若干次内容一分支，将矩阵本身绕回去，剩下的作用就是乘 $g^{-1}$。因此行列式的素因子可以给有限尺度变化。
\InsertProof{cor:finite-group}

\section{从模三走到模九：为什么有时会卡住}
普通平移 $x\mapsto x+1$ 在模 $\ell^b$ 下当然走遍所有类。一般分式变换模 $\ell$ 看起来像平移，却可能在更高精度过早循环。因此要额外检查每走 $\ell$ 步，步长是不是恰好多出一个 $\ell$ 因子。
\InsertProof{09_terminal-unit4}
若 $N^2=cI$，所有高次幂都能写成 $A I+B N$，因为偶次幂给标量，奇次幂给 $N$。于是很多矩阵计算只剩两项。例如
\[(I+N)^3=(1+3c)I+(3+c)N.
\]
当 $c=0\pmod3$，$3+c$ 通常恰含一个三；但 $c=6\pmod9$ 时它含至少两个三，就可能跳过一层。这解释了后面为什么会特意避开一个模九余类。
\InsertProof{lem:local-cycle}

\section{我们到底要把参数送到哪里}
回到 $K_{p,r}$，令 $y=a+pz$，合法性为 $\gcd(y,n)=1$。想选 $j+1$ 份第一块、$j$ 份第二块，以及 $p-1-2j$ 份载体。将它们的和算出来：
\[
(j+1)a+j(rz-a)-(p-1-2j)pz=y-(p^2-nj)z.
\]
所以只要把参数方向送到
\begin{equation}\label{eq:terminal}
t_j=\binom{1}{p^2-nj},\qquad
\max(0,\lceil(p-1-r)/2\rceil)\le j\le(p-1)/2,
\end{equation}
这组数的和就是零。$j$ 的区间只是要求载体取用数在 $0,r$ 之间。

举例：$p=5,n=11,r=1$，区间只允许 $j=2$，终端列是 $(z,y)=(1,3)$。于是 $a=y-pz=-2$，第二块值为 $3$，载体为 $-5$。三份 $-2$ 加两份 $3$ 恰好为零。这是一个非零的配置，但已经能接第5讲的收尾。

\section{为什么只会生成少量小数字，也够用了}
互补混合中，份数最多是 $p$。因此容易得到的是不超过 $p$ 的安全数字，而不是模 $n$ 的每个单位直接给一个操作。下一个引理说，在 $2p<n<4p$ 这段范围，所有小安全数字已经能通过相乘和取逆生成模 $n$ 的全部单位。

证明不是找原根。奇数 $n$ 时，假设有一个最小缺失单位，通过乘二或四再折回较小正代表，会得到更小的缺失单位，矛盾。偶数时则按是否被四整除，利用平衡代表或数量超过群的一半来完成。
\InsertProof{lem:small-units}

\section{把所有工具接成一条有限路径}
下面定理列出三个输入条件，你可以把它看成后面构造的任务单：已经能在同余类内部精确移动；能在模 $n$ 下改变第一坐标；还要能处理有限共同尺度和二进信息。
\InsertProof{thm:terminal}

\section{再慢看一次：模n正确以后，误差怎样被消掉}
最容易被一句“线性修正”带过的是第四步。假设当前列 $w$ 与目标 $t$ 已模 $h$ 相同，写
\[w=t+he.
\]
施加 $I+hX$，模 $h^2$ 下得到
\[(I+hX)w\equiv t+h(e+Xt)\pmod{h^2}.
\]
于是修正整列，只需让 $Xt\equiv-e\pmod h$。如果我们手里的两个矩阵方向作用于 $t$ 后成为两条独立列，这就是一个二元线性方程组。

本题中的两条列为 $(0,1)^T$ 与 $(-2/p^2,2)^T$，行列式是 $2/p^2$。对奇素数部分可逆，所以够用；对二进部分不够，因此才先用乘法调整第一坐标，再用一次下剪切调整第二坐标。这就是定理分开处理奇偶部分的真正原因。

最后列已经与 $t_j$ 模 $h^2$ 完全相同，上一讲的精确纤维引理才出场。这时无限数值大小被处理掉，接下来平均零和组并收尾，全程只有有限次操作。
